(3) is the computation of Theorem \claimref{hawking.i-state}(3) with the stated core in place of the window span: the core is $\tau$-invariant, $e^{-\beta h}$ preserves it and is bounded on it, and density plus boundedness of $A$ and $A_I$ extend the bounded identity $A = A_I$ from the core to $\mathfrak{h}$; the unbounded identity is used only on the core, where $1 + e^{-\beta h}$ acts bijectively. (2) is Theorem \claimref{hawking.i-state}(2) on the same core. (1): the state $\omega_{A_I}$ exists for the symbol $0 \le A_I \le 1$ (Fact \claimref{hawking.imported-quasi-free-states-car-algebra}); its two-point function satisfies the two-point KMS identity by (3) read backwards, and Fact \claimref{hawking.imported-quasi-free-kms-reduces-two} upgrades the two-point identity to the KMS property of the state. What this paper certifies independently, and the script checks exactly, is the window case, where no imported calculus is needed.
